\documentclass[11pt,a4paper]{article}

\usepackage[T1]{fontenc}
\usepackage[utf8]{inputenc}
\usepackage{lmodern}
\usepackage[margin=2.5cm]{geometry}
\usepackage{booktabs}
\usepackage{multirow}
\usepackage{amsmath}
\usepackage{siunitx}
\usepackage{xcolor}
\usepackage[hidelinks]{hyperref}
\usepackage{enumitem}
\setlist{nosep, leftmargin=*}

\newcommand{\ci}[2]{\,[#1,\,#2]}

\title{Does post-treatment MRI add to clinical information for short-term progression prediction in glioblastoma?\\[0.3em]
\large A censoring-aware re-evaluation on MU-Glioma-Post with external testing on LUMIERE}
\author{PRISM project --- internal technical report}
\date{10 October 2026}

\begin{document}
\maketitle

\begin{abstract}
\noindent\textbf{Aim.} To test whether information from a single post-treatment MRI improves prediction of progression
within 120 days over information that is available without imaging (clinical, timing and treatment).
\textbf{Data.} MU-Glioma-Post (University of Missouri; 314 eligible surveillance scans from 157 patients, of which
253 scans / 128 patients are WHO grade~4) for development, and LUMIERE (132 scans / 61 glioblastoma patients) for
external testing.
\textbf{Methods.} Censoring-aware labels; $5\times5$ repeated patient-grouped cross-validation; IPCW AUC at 120 days with
patient-cluster bootstrap confidence intervals; nested selection over feature groups and over 70 radiomics
configurations; paired $\Delta$AUC against a non-imaging baseline.
\textbf{Results.} Demographics and molecular markers alone were weak (AUC 0.56--0.63, grade~4). Adding timing raised
the AUC to 0.758 (gradient boosting). Imaging-only models reached at most 0.660. Adding imaging to the non-imaging
model did not improve it reliably; in the honest nested estimate, radiomics changed the AUC by $-0.008$ (grade~4) and
$+0.013$ (all gliomas), and none of 70 radiomics configurations had a $\Delta$AUC confidence interval above zero.
On the PRISM report split, a model without imaging scored 0.870--0.901 against 0.802 for the PRISM report radiomics
model. On LUMIERE, the PRISM report radiomics model dropped to 0.61, while a clinical + timing + volume model reached
0.72.
\textbf{Conclusion.} In this setting (conventional MRI, one development centre, short-horizon progression), imaging adds
no measurable value over clinical and treatment information.
\end{abstract}

\section{Background}
Surveillance MRI after treatment of glioblastoma is used to decide when to intervene, and several groups have proposed
image-based models that predict progression from a follow-up scan. Reported discrimination on MU-Glioma-Post is
around 0.80, but these figures come from single train/test splits and are rarely compared with a baseline that uses the
clinical and treatment information already available on the day of the scan. The PRISM report radiomics model
(AUC 0.802 on a 30-patient test split) is one such model.

This report asks a narrow question: \emph{once clinical, timing and treatment information is in the model, does the
image add anything?}

\section{Data and endpoint}

\subsection{Development cohort: MU-Glioma-Post}
The TCIA MU-Glioma-Post collection contains 596 post-treatment timepoints from 203 patients, with T1, T1c, T2 and
FLAIR images, expert-refined four-label masks (enhancing tumour, non-enhancing core, surrounding FLAIR hyperintensity,
resection cavity) and a clinical workbook. Eligibility rules were applied in order (Table~\ref{tab:flow}). Index scans
are surveillance scans taken before first progression and before any late treatment (second-line therapy,
immunotherapy, or new treatment).

\begin{table}[h]
\centering
\caption{Cohort flow (MU-Glioma-Post).}
\label{tab:flow}
\begin{tabular}{lrrr}
\toprule
Step & Scans & Patients & Excluded \\
\midrule
All timepoints & 596 & 203 & -- \\
All four sequences and mask present & 594 & 203 & 2 \\
Tumour (labels 1--3) present & 592 & 203 & 2 \\
MRI day known & 589 & 202 & 3 \\
Before late treatment & 508 & 199 & 81 \\
Before first progression (index scans) & \textbf{314} & \textbf{157} & 194 \\
\quad of which WHO grade 4 & 253 & 128 & \\
\bottomrule
\end{tabular}
\end{table}

\subsection{Endpoint}
For each index scan, the label at horizon $H = 120$ days is 1 if progression is documented within $H$ days, 0 if
progression occurs later, the patient died without progression, or progression-free follow-up reaches $H$ days, and
censored otherwise. Follow-up uses the latest dated MRI or clinical event. At 120 days the cohort has 99 positive,
178 negative and 37 censored scans (grade~4: 92 / 136 / 25).

\subsection{External cohort: LUMIERE}
LUMIERE (Suter et al., 2022) provides expert RANO ratings per timepoint for 91 glioblastoma patients. Index scans are
rated follow-ups (stable disease, partial or complete response, post-operative; at least two weeks after surgery)
before the first RANO progression. With the same 120-day rule this gives 132 scans from 61 patients (47 positive,
77 negative, 8 censored). LUMIERE has no treatment fields, so only features available in both cohorts (age, sex,
IDH/MGMT, time since surgery, early post-radiotherapy flag, compartment volumes) are used externally.

\section{Methods}

\subsection{Feature groups}
\begin{itemize}
\item \textbf{Clinical}: age, sex, IDH and MGMT status with an explicit ``unknown'' level.
\item \textbf{Timing}: days since diagnosis, days since end of radiotherapy, early post-radiotherapy flag
(within 12 weeks), visit index.
\item \textbf{Treatment}: radiotherapy given, dose, fractions, hypofractionation, start day and duration; initial
chemotherapy; day of surgery; biopsy before resection; previous brain tumour.
\item \textbf{Volumes}: log volumes of each mask compartment.
\item \textbf{Morphology}: enhancing fraction, multifocality, largest axial enhancing area, distance of enhancement to
the cavity, intensities relative to normal-appearing brain.
\item \textbf{Change}: differences from the previous scan and from the nadir, growth rates, a $\geq 40\%$ enhancing-volume
increase flag, new enhancing component.
\item \textbf{Radiomics}: the PRISM report features (two grey levels) and corrected PyRadiomics features (whole-brain
normalisation, 32 bins, per compartment), with 70 configurations of stability filtering, feature selection and
integration (concatenation or a stacked radiomics score).
\end{itemize}

\subsection{Models and evaluation}
Two learners were used: ridge logistic regression with the penalty chosen by inner patient-grouped cross-validation, and
shallow gradient boosting. All preprocessing and selection were fitted within training folds. Models were evaluated with
$5\times5$ repeated patient-grouped stratified cross-validation, using inverse-probability-of-censoring-weighted (IPCW)
cumulative/dynamic AUC at 120 days. Confidence intervals come from a patient-cluster bootstrap (1000 resamples).
Paired $\Delta$AUC was computed against the ridge non-imaging baseline. Nested selection over feature groups and over
radiomics configurations gives an estimate that includes the cost of choosing a configuration. A label-permutation
control returned AUC $\approx 0.50$.

\section{Results}

\subsection{Imaging versus non-imaging information}
Table~\ref{tab:main} compares models by the information they use.

\begin{table}[h]
\centering
\caption{IPCW AUC at 120 days, MU-Glioma-Post grade 4 (253 scans, 128 patients), repeated patient-grouped
cross-validation, 95\% patient-cluster bootstrap CI.}
\label{tab:main}
\small
\begin{tabular}{lll}
\toprule
Information used & Ridge & Boosting \\
\midrule
\multicolumn{3}{l}{\emph{Without imaging}} \\
Clinical only (age, sex, IDH/MGMT) & 0.556\ci{0.456}{0.647} & 0.634\ci{0.545}{0.713} \\
Clinical + timing & 0.662\ci{0.583}{0.754} & \textbf{0.758}\ci{0.677}{0.838} \\
Clinical + timing + treatment & 0.694\ci{0.613}{0.775} & 0.723\ci{0.633}{0.806}$^{a}$ \\
\midrule
\multicolumn{3}{l}{\emph{Imaging only}} \\
Volumes + morphology + change & 0.603\ci{0.515}{0.686} & 0.545\ci{0.462}{0.628} \\
\quad + PRISM report radiomics & 0.660\ci{0.589}{0.736} & 0.652\ci{0.567}{0.730} \\
\midrule
\multicolumn{3}{l}{\emph{Non-imaging + imaging}} \\
Clinical + timing + volumes & 0.681\ci{0.604}{0.768} & 0.733\ci{0.660}{0.805} \\
Clinical + timing + radiomics & 0.709\ci{0.631}{0.783} & 0.740\ci{0.669}{0.816} \\
Clinical + timing + all imaging & 0.687\ci{0.610}{0.767} & 0.717\ci{0.643}{0.797} \\
\bottomrule
\end{tabular}

\vspace{0.3em}
\footnotesize $^{a}$ Regularised boosting (depth 2); other boosting rows use depth 3 and overfit more
(training AUC 0.96--1.00).
\end{table}

All gliomas (314 scans, 157 patients) show the same pattern: clinical + timing with boosting reaches 0.777; adding
radiomics gives 0.777, volumes 0.760 and all imaging 0.752; imaging-only models reach at most 0.695.

\subsection{Does radiomics add to the non-imaging baseline?}
Table~\ref{tab:c3} gives the paired comparison against the clinical + timing + treatment baseline.

\begin{table}[h]
\centering
\caption{Radiomics added to the clinical + timing + treatment baseline (ridge). $\Delta$AUC is paired against the
baseline.}
\label{tab:c3}
\footnotesize
\begin{tabular}{lllll}
\toprule
& \multicolumn{2}{c}{Grade 4} & \multicolumn{2}{c}{All gliomas} \\
\cmidrule(lr){2-3}\cmidrule(lr){4-5}
Model & AUC & $\Delta$AUC & AUC & $\Delta$AUC \\
\midrule
Baseline & 0.694 & -- & 0.719 & -- \\
Best of 70 configurations (post hoc) & 0.720 & $+0.024$\ci{-0.032}{+0.079} & 0.754 & $+0.034$\ci{-0.018}{+0.082} \\
Nested choice over 70 (honest) & 0.687 & $-0.008$\ci{-0.046}{+0.026} & 0.733 & $+0.013$\ci{-0.031}{+0.054} \\
\midrule
Configurations with $\Delta$AUC CI $> 0$ & \multicolumn{2}{c}{0 / 70} & \multicolumn{2}{c}{0 / 70} \\
\bottomrule
\end{tabular}
\end{table}

Correcting the radiomics extraction did not help. The small signal in the PRISM report features comes from texture
computed on two grey levels, which acts as a description of the spatial pattern of bright and dark regions within the
lesion; the corrected 32-bin features carried no additional signal.

In nested selection over feature groups, clinical + timing was part of the chosen combination in 22 of 25 outer folds
(grade~4) and 24 of 25 (all gliomas). An imaging group was usually added as well, but a different one from fold to
fold, which is consistent with imaging contributing noise-level rather than stable signal.

\subsection{PRISM report split}
On the 30 test patients of the PRISM report split, with configurations fixed beforehand from cross-validation
(Table~\ref{tab:split}), a model without imaging outperformed the PRISM report radiomics model, and adding radiomics
changed it by at most $+0.005$. The same models average 0.72--0.74 under repeated cross-validation, so this split is
favourable to every model.

\begin{table}[h]
\centering
\caption{PRISM report split (127 training / 30 test patients, all gliomas).}
\label{tab:split}
\small
\begin{tabular}{lcc}
\toprule
Model & AUC (original label) & IPCW AUC (120 d) \\
\midrule
PRISM report radiomics model & 0.802 & 0.811 \\
Clinical + timing + treatment, ridge & 0.870 & 0.844 \\
Clinical + timing + treatment, boosting & \textbf{0.901} & \textbf{0.881} \\
\quad + PRISM report radiomics (top 10) & 0.869 & 0.848 \\
\quad + corrected compartment radiomics (stacked) & 0.869 & 0.843 \\
\bottomrule
\end{tabular}
\end{table}

\subsection{External test on LUMIERE}
Models trained on all MU grade-4 scans were applied unchanged to LUMIERE (Table~\ref{tab:ext}). The PRISM report
radiomics model, re-extracted with its original pipeline on 118 of the 132 scans, fell to 0.613. Among shared-feature
models, volumes helped the boosting model (0.663 to 0.719) but not the ridge model (0.628 to 0.629); the confidence
intervals overlap. Recalibrating the intercept and slope on other LUMIERE patients brought the observed/expected ratio
to 0.94--0.98 without changing the ranking.

\begin{table}[h]
\centering
\caption{Locked models on LUMIERE (132 scans, 61 patients), pooled IPCW AUC at 120 days.}
\label{tab:ext}
\small
\begin{tabular}{llcc}
\toprule
Features & Learner & MU CV & LUMIERE [95\% CI] \\
\midrule
PRISM report radiomics model & -- & -- & 0.613\ci{0.500}{0.729} \\
Clinical & ridge & 0.578 & 0.609\ci{0.479}{0.735} \\
Clinical + timing & ridge & 0.644 & 0.628\ci{0.521}{0.748} \\
Clinical + timing & boosting & 0.714 & 0.663\ci{0.513}{0.781} \\
Volumes & ridge & 0.591 & 0.534\ci{0.427}{0.629} \\
Clinical + timing + volumes & ridge & 0.695 & 0.629\ci{0.505}{0.753} \\
Clinical + timing + volumes & boosting & 0.709 & \textbf{0.719}\ci{0.609}{0.832} \\
\bottomrule
\end{tabular}
\end{table}

\section{Limitations}
\begin{itemize}
\item \textbf{Treatment features are not yet restricted to what was known at the scan date.} Radiotherapy dose, duration
and chemotherapy are used even for scans taken before radiotherapy ended (29 of 215 such scans). A one-off check with
as-of-scan masking gave a \emph{higher} baseline (grade~4 ridge 0.694 to 0.732), so the conclusion is not driven by this,
but the tables will be re-run once the masking is in the code.
\item \textbf{The label is partly shaped by the scan schedule.} Progression is usually declared at an MRI: 61 of 106
progressing patients have their progression date within three days of an MRI, and the interval to the next MRI alone
(not available at prediction time) gives AUC 0.64 on grade~4. Part of the timing signal may reflect closer follow-up
early after radiotherapy rather than tumour biology.
\item \textbf{Scope.} Conventional MRI only (no perfusion or diffusion), a single development centre, about 130--160
patients, and a 120-day horizon. The result does not show that imaging is uninformative about progression in general.
\item \textbf{Multiple comparisons.} Many configurations were evaluated on the same folds. Post-hoc rows are optimistic;
the nested and paired estimates are the ones to quote. The LUMIERE result of 0.719 is the best of 16 combinations.
\item \textbf{External cohort differences.} LUMIERE defines progression by expert RANO rating on imaging, uses
automatic segmentation in atlas space (its necrosis compartment likely includes the resection cavity), and lacks
treatment fields.
\end{itemize}

\section{Conclusion}
A model that uses no imaging --- clinical, timing and treatment information --- performs as well as or better than models
that also use the MRI, and imaging alone is weaker. No imaging addition reached statistical significance on
MU-Glioma-Post. The supported claim is that, in this setting, \emph{imaging does not add measurable value over clinical
and treatment information}; it is not that clinical information beats imaging in general. Image-based progression
models should be reported against a treatment-aware non-imaging baseline under patient-grouped, censoring-aware
evaluation.

\section*{Data and code}
Pipeline: \texttt{prism/} in this repository. Source results: the C2, C3, external LUMIERE and imaging-versus-non-imaging
write-ups in \texttt{docs/results/}.

\begin{thebibliography}{9}
\bibitem{mu} MU-Glioma-Post. The Cancer Imaging Archive, \url{https://doi.org/10.7937/7K9K-3C83}; data descriptor Sci Data,
\url{https://doi.org/10.1038/s41597-025-06011-7}.
\bibitem{lumiere} Suter Y, et al. The LUMIERE dataset: longitudinal glioblastoma MRI with expert RANO evaluation.
Sci Data 2022, \url{https://doi.org/10.1038/s41597-022-01881-7}.
\end{thebibliography}

\end{document}
